% 来源：docs/07-generative-adversarial-networks.md
% 仅建立章节框架；正文、公式、图表与习题内容待后续迁移。
\chapter{生成对抗网络}
\label{chap:gan}

\section*{本章要点}
% TODO：迁移原章“先抓住这几件事”的要点。
\contentplaceholder

\section{为什么考虑不依赖似然的训练？}
\label{sec:gan-motivation}
% TODO：补充本节内容。
\contentplaceholder

\subsection{似然、视觉质量与样本多样性}
\label{subsec:gan-evaluation}
% TODO：补充本节内容。
\contentplaceholder

\subsection{双样本比较与分类器检验}
\label{subsec:gan-two-sample}
% TODO：补充本节内容。
\contentplaceholder

\section{两个网络与原始目标}
\label{sec:gan-objective}
% TODO：补充本节内容。
\contentplaceholder

\subsection{生成器与隐式分布}
\label{subsec:gan-generator}
% TODO：补充本节内容。
\contentplaceholder

\subsection{判别器与真假分类}
\label{subsec:gan-discriminator}
% TODO：补充本节内容。
\contentplaceholder

\subsection{原始 minimax 对抗目标}
\label{subsec:gan-minimax}
% TODO：补充本节内容。
\contentplaceholder

\section{最优判别器与 JS 散度}
\label{sec:gan-js}
% TODO：补充本节内容。
\contentplaceholder

\subsection{固定生成器时的最优判别器}
\label{subsec:gan-optimal-discriminator}
% TODO：补充本节内容。
\contentplaceholder

\subsection{JS 散度的定义与性质}
\label{subsec:gan-js-definition}
% TODO：补充本节内容。
\contentplaceholder

\subsection{最优判别器下目标的推导}
\label{subsec:gan-derivation}
% TODO：补充本节内容。
\contentplaceholder

\subsection{分布匹配的平衡点与理论条件}
\label{subsec:gan-equilibrium}
% TODO：补充本节内容。
\contentplaceholder

\section{实际训练：交替更新与非饱和损失}
\label{sec:gan-training}
% TODO：补充本节内容。
\contentplaceholder

\subsection{判别器与生成器的交替更新}
\label{subsec:gan-alternating}
% TODO：补充本节内容。
\contentplaceholder

\subsection{梯度饱和与非饱和生成器损失}
\label{subsec:gan-non-saturating}
% TODO：补充本节内容。
\contentplaceholder

\subsection{模式坍塌与训练不稳定}
\label{subsec:gan-collapse}
% TODO：补充本节内容。
\contentplaceholder

\section{WGAN：从分类差异到搬运距离}
\label{sec:gan-wgan}
% TODO：补充本节内容。
\contentplaceholder

\subsection{为什么 JS 可能缺乏距离信息？}
\label{subsec:gan-support}
% TODO：补充本节内容。
\contentplaceholder

\subsection{最优传输与对偶形式}
\label{subsec:gan-transport}
% TODO：补充本节内容。
\contentplaceholder

\subsection{Critic 的目标与 Lipschitz 约束}
\label{subsec:gan-critic}
% TODO：补充本节内容。
\contentplaceholder

\subsection{权重裁剪及其局限}
\label{subsec:gan-clipping}
% TODO：补充本节内容。
\contentplaceholder

\subsection{WGAN-GP 的梯度惩罚}
\label{subsec:gan-gradient-penalty}
% TODO：补充本节内容。
\contentplaceholder

\section{对抗损失可以与其他任务组合}
\label{sec:gan-applications}
% TODO：补充本节内容。
\contentplaceholder

\subsection{条件 GAN 与复合损失}
\label{subsec:gan-conditional}
% TODO：补充本节内容。
\contentplaceholder

\subsection{SRGAN、图像修复与 VQ-GAN}
\label{subsec:gan-restoration}
% TODO：补充本节内容。
\contentplaceholder

\subsection{Pix2Pix：成对图像翻译}
\label{subsec:gan-pix2pix}
% TODO：补充本节内容。
\contentplaceholder

\subsection{CycleGAN：非成对翻译与循环一致性}
\label{subsec:gan-cyclegan}
% TODO：补充本节内容。
\contentplaceholder

\subsection{StarGAN 与 ADDA：多域建模与域适应}
\label{subsec:gan-domains}
% TODO：补充本节内容。
\contentplaceholder

\section{从高分辨率生成到可控生成}
\label{sec:gan-controllability}
% TODO：补充本节内容。
\contentplaceholder

\subsection{Progressive Growing GAN}
\label{subsec:gan-progressive}
% TODO：补充本节内容。
\contentplaceholder

\subsection{StyleGAN 的中间潜在空间与样式控制}
\label{subsec:gan-stylegan}
% TODO：补充本节内容。
\contentplaceholder

\subsection{StyleGAN2 的结构改进}
\label{subsec:gan-stylegan2}
% TODO：补充本节内容。
\contentplaceholder

\subsection{StyleGAN3 与后续阅读}
\label{subsec:gan-future}
% TODO：补充本节内容。
\contentplaceholder

\section{自检}
\label{sec:gan-review}
% TODO：补充本节内容。
\contentplaceholder

\section{来源与延伸阅读}
\label{sec:gan-reading}
% TODO：补充本节内容。
\contentplaceholder
